\section{Les outils de développement}

%-------------------------------------------------------------------------------
%
%-------------------------------------------------------------------------------
\subsection{La compilation}

Observons les principaux éléments du {\tt Makefile}.

%...............................................................................
%
%...............................................................................
\subsubsection{L'assembleur}

   J'utilise NASM pour les quelques fichiers écrits en assembleur.

%...............................................................................
%
%...............................................................................
\subsubsection{Le compilateur C}

%...............................................................................
%
%...............................................................................
\subsubsection{La création du noyau}

   Sans grande surprise, on compile avec

\begin{lstlisting}{}
  make
\end{lstlisting}

  Différentes cibles sont créées

\begin{description}
   \item [{\tt manux}]
   \item [{\tt run}]
   \item [{\tt clean}]
\end{description}

%-------------------------------------------------------------------------------
%
%-------------------------------------------------------------------------------
\subsection{L'exécution}

   Le but n'est évidemment pas de faire tourner cette merveille sur du
vrai matériel, mais plutôt dans un émulateur.

%...............................................................................
%
%...............................................................................
\subsubsection{L'émulateur Bochs}

   J'utilisais initialement cet outil, mais les fichiers de
configuration ne semblent plus convenir aux dernières versions de
Bochs, il faudrait les mettre à jour ! Des volontaires ?

%...............................................................................
%
%...............................................................................
\subsubsection{L'émulateur QEMU}

   C'est maintenant {\tt qemu} que j'utilise, par exemple de la façon
suivante
   
\begin{lstlisting}{}
qemu-system-i386 -drive format=raw,file=manux,index=0,if=floppy -m 64M
\end{lstlisting}

  C'est d'ailleurs la commande que lance

\begin{lstlisting}{}
   make run
\end{lstlisting}

%-------------------------------------------------------------------------------
%
%-------------------------------------------------------------------------------
\subsection{Le débugage ``interne''}

   Ce que j'appelle ici le débugage interne, c'est tous les outils que
l'on peut utiliser dans le code lui-même pour essayer de comprendre ce
qu'il s'y passe.

   Nous avons pour le moment quelques outils de base, copieusement bugués
eux-mêmes (!).

%...............................................................................
%
%...............................................................................
\subsubsection{La fonction {\tt printk\_debug}}

   Elle permet de produire des messages d'erreur qui seront envoyés de
façon sélective en fonction de la valeur de \lstinline!masqueDebugage!
qui doit être construite dans {\tt manux/debug.h}.

%...............................................................................
%
%...............................................................................
\subsubsection{La macro {\tt assert}}

   Voir {\tt manux/debug.h}

%...............................................................................
%
%...............................................................................
\subsubsection{Le handler {\tt handlerPanique}}
 
%...............................................................................
%
%...............................................................................
\subsubsection{Le fichier {\tt noyau/symbols}}

   Il va surtout nous servir pour aider au débugage externe, mais il
est généré à la compilation du noyau.
 
%-------------------------------------------------------------------------------
%
%-------------------------------------------------------------------------------
\subsection{Le débugage ``externe''}

   Une technique pour débuguer est d'utiliser {\tt qemu} couplé avec
{\tt gdb}. Pour cela, je lance par exemple {\tt qemu} de la façon suivante

\begin{lstlisting}{}
qemu-system-i386 -drive format=raw,file=manux,index=0,if=floppy  -rtc
base=localtime -m 64M -gdb tcp::1234  -S 
\end{lstlisting}

   Je lance ensuite {\tt gdb} dans un autre terminal avec un
\lstinline!.gdbinit! qui contient :

\begin{lstlisting}{}
  target remote :1234
  dir noyau
  dir lib
  file noyau/noyau.elf
\end{lstlisting}

   La commande \lstinline!dir! permet de dire à {\tt gdb} où trouver
les fichiers sources et la command \lstinline:file: lui permet de
consulter le binaire du noyau pour y trouver l'adresse des fonctions
et variables (ici en fait, c'est une autre compilation du noyau au
format {\sc elf} !).

%...............................................................................
%
%...............................................................................
\subsubsection{Exécution}

   Quelques commandes de base pour l'exécution de code dans {\tt gdb}

\begin{description}
     \item[{\tt {\bf r} un}] permet de lancer l'exécution d'un
       programme (inutilisable ici étant donné notre fonctionnement :
       on débogue à distance) ;
     \item[{\tt {\bf c} ontinue}] permet de poursuivre l'exécution jusqu'au
       prochain point d'arrêt/bug/fin/\ldots ;
     \item[{\tt {\bf s} tep}] avance d'une instruction dans le code ;
     \item[{\tt {\bf n} ext}]
     \item[{\tt {\bf b} reakpoint fonction}]
     \item[{\tt {\bf b} reakpoint fichier:ligne}]
\end{description}

%...............................................................................
%
%...............................................................................
\subsubsection{Affichage des registres}

   Il peut être intéressant de consulter le contenu des registres

\begin{description}
     \item[{\tt info registers}] donne le contenu des principaux registres.
\end{description}

%...............................................................................
%
%...............................................................................
\subsubsection{Affichage des variables}

   {\tt Gdb} permet également de consulter le contenu des variables :

\begin{description}
   \item[{\tt p toto}] affiche le contenu de la variable {\tt toto}.
\end{description}

%...............................................................................
%
%...............................................................................
\subsubsection{Désassemblage du code}

    Une fonctionnalité intéressante est de pouvoir observer le code
assembleur, ce qui peut permettre de comprendre certains comportements
étranges, tels que les conséquences de choix d'optimisation du
compilateur :

\begin{description}
   \item[{\tt disassemble fonc}] affiche le code assembleur de la
     fonction {\tt fonc}.
\end{description}

